\section{Ordinary backtracking and a computable certificate}
\label{sec:model}
Let $Q\subset\R^d$ be nonempty, compact and convex, with diameter
$D>0$. Let $f$ be convex and $C^2$ on a neighborhood of $Q$,
$f_Q^*=\min_Q f$, and $\delta(x)=f(x)-f_Q^*$. Assume
\begin{equation}\label{eq:curvature}
 \norm{\nabla^2f(z)}_2\le H_0+H_1\delta(z),\qquad z\in Q,
 \quad H_0,H_1\ge0.
\end{equation}
The reference value in this assumption is the constrained optimum.
A bound relative to a smaller unconstrained infimum does not imply the
same constrained bound without modifying $H_0$.
If $D=0$ the only feasible point solves the problem.

At a source $x_k$, obtain an exact linear minimization oracle (LMO),
and form the FW gap and accumulated certificate:
\begin{align}
 s_k&\in\arg\min_{v\in Q}\inner{\nabla f(x_k)}{v},\quad
 G_k=\inner{\nabla f(x_k)}{x_k-s_k},\label{eq:gap}\\
 \ell_k&=\max_{0\le i\le k}\{f(x_i)-G_i\},\quad
 U_k=f(x_k)-\ell_k.\label{eq:certificate}
\end{align}
Convexity gives $0\le\delta_k\le U_k\le G_k$. A valid known lower
bound can be included in $\ell_k$, but is not required. A zero gap
certifies optimality. We stop at $U_k\le\varepsilon$.

\paragraph{The algorithm.}
Choose $M>0$ and set $L_{-1}=M$. At each nonterminal source, put
$d_k=s_k-x_k$, $q_k=\norm{d_k}$, and start with $L=L_{k-1}/2$.
Repeat
\begin{equation}\label{eq:btstep}
 \alpha=\min\{1,G_k/(Lq_k^2)\},\qquad
 f(x_k+\alpha d_k)\le f(x_k)-\alpha G_k
                         +\tfrac12L\alpha^2q_k^2
\end{equation}
until the inequality passes, doubling $L$ after each rejection.
Store $L_k=L$, update $x_{k+1}=x_k+\alpha d_k$, and reuse the accepted
value. Only trial function values are recomputed during the search.
This is the standard quadratic backtracking scheme~\cite{pedregosa}.
For $M$ fixed, the algorithm does not need $D,H_0,H_1$, or $f_Q^*$.
The optional $M=U_0/D^2$ uses a diameter bound and the first paid
function, gradient and LMO calls.

\begin{lemma}[Local acceptance]\label{lem:local}
For a feasible descent direction $d$ with $g=-\inner{\nabla f(x)}d>0$
and feasible cap $a_{\max}>0$, the test in~\eqref{eq:btstep} passes
at $\alpha=\min\{a_{\max},g/(L\norm d^2)\}$ whenever
$L\ge H_0+H_1\delta(x)$. Every accepted trial, regardless of this
sufficient threshold, satisfies
\begin{equation}\label{eq:acceptedprogress}
 f(x)-f(x+\alpha d)\ge\tfrac12\alpha g.
\end{equation}
Thus the search terminates and
\begin{equation}\label{eq:localenvelope}
 L_k\le\max\{L_{k-1}/2,\ 2(H_0+H_1\delta_k)\}.
\end{equation}
The assertion also allows $a_{\max}>1$.
\end{lemma}
\begin{proof}
Write $h(t)=f(x+td)-f(x)$ and $K=H_0+H_1\delta(x)$.
While $h\le0$, the Hessian bound gives $h''(t)\le K\norm d^2$.
Initially $h'(0)=-g<0$. If $K>0$, a first return to zero before
$\min\{a_{\max},2g/(K\norm d^2)\}$ contradicts
$h(t)\le-gt+K\norm d^2t^2/2<0$. The same integration proves
the quadratic upper bound through the prescribed trial; the zero-$K$
case follows with a linear upper model. Its curvature is at most $L$,
so the test passes. For any accepted trial,
$L\alpha\norm d^2\le g$, proving~\eqref{eq:acceptedprogress}.
If there is a rejection, the last rejected estimate is less than $K$,
so the accepted one is less than $2K$; otherwise it is $L_{k-1}/2$.
\end{proof}

Consequently $f_k$ and $U_k$ are nonincreasing. The analysis uses
curvature only on a descent segment, not a uniform upper model on
every segment point.

\begin{lemma}[Lag potential]\label{lem:lag}
Let source phase indices be nondecreasing in $\{0,\ldots,J-1\}$.
Let positive caps obey $C_j/2\le C_{j+1}\le C_j$.
Suppose an update with $L_k>C_j$ in its source phase must have
$L_k=L_{k-1}/2$. Then the number of such lag updates is at most
\[
 \ceil{\log_2\max\{1,M/C_0\}}+J-1.
\]
The statement applies to every prefix and allows initially skipped phases.
\end{lemma}
\begin{proof}
Before an update in phase $j$, use the nonnegative integer potential
\[
 p=\max\{0,\ceil{\log_2(L_{k-1}/C_j)}\}.
\]
A lag update reduces it by one; a non-lag update leaves zero potential
before a phase change. A forward jump of $m$ phases raises it by at
most $m$. Starting from virtual phase zero, all jumps, including an
initial jump, sum to at most $J-1$. Nonnegativity proves the count.
\end{proof}

\begin{theorem}[Observable stopping and total function cost]\label{thm:fw}
For $0<\varepsilon<U_0$, let
\[
 J=\ceil{\log_2(U_0/\varepsilon)},\qquad
 P=\ceil{\log_2\max\{1,MD^2/U_0\}}.
\]
Ordinary backtracking reaches $U_N\le\varepsilon$ within
\begin{equation}\label{eq:fwcount}
 N\le P+\ceil{(6+16H_1D^2)J+32H_0D^2/\varepsilon}
\end{equation}
updates. It uses $N+1$ gradients and LMOs, including the final check.
For its actual $N\ge1$, every function evaluation, including all
rejected trials, is counted in
\begin{equation}\label{eq:functioncount}
 N_f\le1+2N+\ceil{\log_2\max\{1,2(H_0+H_1U_0)/M\}}.
\end{equation}
For $N=0$, one function, gradient and LMO evaluation suffice.
These are exact-oracle and exact-arithmetic bounds.
\end{theorem}
\begin{proof}
Put $b_j=U_0/2^j$ and assign a nonterminal source to phase $j$ when
$b_j/2<U_k\le b_j$, for $j=0,\ldots,J-1$.
Use
\[
 C_j=\max\{2(H_0+H_1b_j),\ b_j/(2D^2)\}.
\]
The caps decrease by at most a factor two. By
Lemma~\ref{lem:local}, an accepted estimate above $C_j$ must be a
first-trial halving. Since $C_0\ge U_0/(2D^2)$,
Lemma~\ref{lem:lag} bounds lag updates by $P+J$.

For a non-lag update in phase $j$, $G_k\ge U_k>b_j/2$ and
$q_k\le D$ imply
\[
 f_k-f_{k+1}\ge p_j:=\min\{b_j/4,\ b_j^2/(8C_jD^2)\}.
\]
Every update decreases $f$, and the total decrease within the phase
is at most its initial suboptimality, hence at most $b_j$.
Thus its non-lag updates number at most
\[
 \lfloor b_j/p_j\rfloor+1
 \le5+16H_1D^2+16H_0D^2/b_j.
\]
Summing and using $\sum_{j<J}1/b_j<2/\varepsilon$ proves
\eqref{eq:fwcount}. This bounds every unfinished prefix, which also
establishes termination.

The local envelope and monotone $U$ give
$L_k\le\max\{M/2^{k+1},2(H_0+H_1U_0)\}$.
If $r_k$ denotes rejections, $L_k=2^{r_k-1}L_{k-1}$, so exactly
\[
 \sum_{k<N}(r_k+1)=2N+\log_2(L_{N-1}/M)
\]
trial values are computed. Add the initial value and use the envelope
to obtain~\eqref{eq:functioncount}. Accepted values are cached.
\end{proof}

The last identity is standard~\cite{pedregosa}. For $M=U_0/D^2$,
$P=0$, and the function-count overhead is independent of
$\varepsilon$. This normalization is optional for the order of the bound.

\begin{proposition}[Initialization is visible even for affine objectives]
\label{prop:init}
Let $Q=[-1,1]$, $f(x)=-cx$, $c>0$, $x_0=-1$, and $L_{-1}=M>0$.
For $0<\varepsilon<2c$, halving backtracking uses exactly
\[
 N=\ceil{\log_2\left(1+\frac Mc(1-\varepsilon/(2c))\right)}
\]
updates. Thus halving $U_0$ may require $\Theta(\log(M/c))$ updates.
\end{proposition}
\begin{proof}
Every test accepts its first trial: $L_k=M/2^{k+1}$,
$x_{k+1}=\min\{1,x_k+c/L_k\}$. The tangent lower bound is always
$-c$, so $x_N=\min\{1,-1+(2c/M)(2^N-1)\}$ and $U_N=c(1-x_N)$.
Solving the stopping inequality gives the formula.
\end{proof}
This lower bound concerns the specified halving rule, not arbitrary
first-order algorithms. With $M=U_0/D^2=c/2$ the first step solves it.

\paragraph{Comparison with a global bound.}
For $f(x)=\cosh(x)-1$ on $[-R,R]$, $H_0=H_1=1$, while a global
gradient-Lipschitz constant is $L_Q=\cosh R$. Starting at $R$ gives
$U_0=2R\sinh R$. For fixed tolerance and fixed $M$, our update bound
grows as $O(R^3+R^2/\varepsilon)$ as $R\to\infty$, rather than the
exponential $L_QD^2/\varepsilon$ term obtained by substituting
$H_0=L_Q,H_1=0$ in the same certificate theorem.
This compares upper bounds, not a runtime or oracle lower bound;
exact segment minimization solves this particular one-dimensional example
in one step. Appendix~\ref{app:explicit} gives further step comparisons.
